\documentclass[11pt]{amsart}
\usepackage{amsmath,amssymb,amsthm,booktabs}
\usepackage[margin=1.15in]{geometry}
\usepackage{hyperref}
\newtheorem{theorem}{Theorem}
\newtheorem{proposition}[theorem]{Proposition}
\newtheorem{lemma}[theorem]{Lemma}
\newtheorem{corollary}[theorem]{Corollary}
\newtheorem{conjecture}[theorem]{Conjecture}
\theoremstyle{remark}
\newtheorem{remark}[theorem]{Remark}
\newcommand{\PP}{\mathbb{P}}
\newcommand{\EE}{\mathbb{E}}
\newcommand{\cF}{\mathcal{F}}
\newcommand{\one}{\mathbf{1}}
\newcommand{\logit}{\operatorname{logit}}

\title{A non-geodesic edge that survives in the loop-erased ant process}
\author{Yitzchak Shmalo}
\date{November 2026}

\begin{document}
\begin{abstract}
In the loop-erased ant process of Kious, Mailler and Schapira, ants walk one after another from a nest $N$ to a food
source $F$ on a finite graph, each ant stepping along an edge with probability proportional to its weight and then
adding one to the weight of every edge of the loop-erasure of its time-reversed trajectory. It was conjectured that the
weight of an edge grows linearly exactly when the edge lies on a geodesic from $N$ to $F$. Let $G_\ell$ be the graph
made of a triangle $N,u,v$ and two paths of $\ell$ edges, one from $u$ to $F$ and one from $v$ to $F$. We prove that on
$G_\ell$ the normalised weights converge almost surely. The limit is $\rho$ on the edges of the path
$N,u,\dots,F$, it is $1-\rho$ on the edges of the path $N,v,\dots,F$, and it is $(1-1/\ell)\rho(1-\rho)$ on the edge
$uv$, where $\rho$ is a non-degenerate random variable with values in $(0,1)$. For $\ell\ge2$ the edge $uv$ lies on no
geodesic and keeps a positive share of the weight, so the conjecture fails on $G_\ell$, and the proportion of ants that
reinforce a path which is not a geodesic has a positive limit. For $\ell=1$ the graph is the losange of Kious, Mailler
and Schapira, the edge $uv$ dies, and the theorem gives for the loop-erased rule the convergence that they proved for
the uniform-geodesic rule.
\end{abstract}
\maketitle

\section{Introduction}

\subsection{The process and the conjecture}
Let $G=(V,E)$ be a finite connected graph with two distinct marked vertices $N$ (the nest) and $F$ (the food). The
\emph{loop-erased ant process} of Kious, Mailler and Schapira~\cite{KMS} is a sequence of weight vectors
$W(n)=(W_e(n))_{e\in E}$, $n\ge0$, defined as follows. Initially $W_e(0)=1$ for every edge. Given $W(n)$, the
$(n+1)$-st ant performs a random walk $X=(X_0,\dots,X_K)$ on $G$ started at $X_0=N$ and stopped at the first time $K$
at which it reaches $F$; from a vertex $x$ it moves to a neighbour $y$ with probability proportional to
$W_{\{x,y\}}(n)$. Let $\bar X_j=X_{K-j}$ be the reversed trajectory, and let $\gamma_{n+1}$ be its loop-erasure: the
path $\gamma(0),\gamma(1),\dots$ with $\gamma(i)=\bar X_{j_i}$, where $j_0=0$ and
$j_{i+1}=1+\max\{j:\bar X_j=\bar X_{j_i}\}$, stopped when it reaches $N$. This is a simple path from $F$ to $N$, and
\[
W_e(n+1)=W_e(n)+\one\{e\in\gamma_{n+1}\},\qquad e\in E .
\]

\begin{conjecture}[{\cite[Conjecture~1.1(1)]{KMS}}]\label{conj}
For the loop-erased ant process on any finite graph, almost surely $W_e(n)/n\to\chi_e$ for all $e\in E$, where
$\chi_e\ne0$ almost surely if and only if the edge $e$ belongs to at least one of the geodesics from $N$ to $F$.
\end{conjecture}

Kious, Mailler and Schapira prove the conjecture for series-parallel graphs whose source and sink are $N$ and
$F$~\cite[Theorem~1.3]{KMS}. For the simplest graph outside this class, the losange with vertices $N,u,v,F$ and edges
$Nu$, $Nv$, $uv$, $uF$, $vF$, they prove the corresponding statement for a second reinforcement rule, the
uniform-geodesic rule~\cite[Theorem~1.4]{KMS}, and they write that for the loop-erased rule they could not prove
convergence on the losange, although they believe that the same result holds. The general case is the first open
problem of the survey~\cite{Aletti}, where the conjecture is stated in the form that the only edges whose weight
grows linearly are those lying on a geodesic between $N$ and $F$. A variant of the process with two nests is studied
in~\cite{MV}, and the process in which every edge crossed by the walk is reinforced in~\cite{KMS2}.

\subsection{The result}
For $\ell\ge1$ let $G_\ell$ be the graph with vertices $N,u,v,F$, $x_1,\dots,x_{\ell-1}$ and $y_1,\dots,y_{\ell-1}$,
whose edges are $Nu$, $Nv$, $uv$, the edges of the path $u=x_0,x_1,\dots,x_\ell=F$ (the \emph{$u$-leg}) and the edges
of the path $v=y_0,y_1,\dots,y_\ell=F$ (the \emph{$v$-leg}). We call $uv$ the \emph{rung}. The graph $G_1$ is the
losange. The distance from $N$ to $F$ is $\ell+1$, and there are two geodesics, $N,u,x_1,\dots,F$ and
$N,v,y_1,\dots,F$. A path from $N$ to $F$ through the rung has $\ell+2$ edges, so the rung lies on no geodesic.

\begin{theorem}\label{thm:main}
Let $\ell\ge1$ and consider the loop-erased ant process on $G_\ell$. There is a random variable $\rho$ with values in
$(0,1)$ such that, almost surely as $n\to\infty$,
\[
\frac{W_e(n)}{n}\longrightarrow
\begin{cases}
\rho & \text{if $e=Nu$ or $e$ is an edge of the $u$-leg,}\\
1-\rho & \text{if $e=Nv$ or $e$ is an edge of the $v$-leg,}\\
\bigl(1-\tfrac1\ell\bigr)\rho(1-\rho) & \text{if $e=uv$.}
\end{cases}
\]
The variables $\rho$ and $1-\rho$ have the same law, and $\rho$ is not almost surely constant.
\end{theorem}

\begin{corollary}\label{cor:fail}
Let $\ell\ge2$. On $G_\ell$ the rung lies on no geodesic from $N$ to $F$ and $\chi_{uv}=(1-1/\ell)\rho(1-\rho)>0$
almost surely, so Conjecture~\ref{conj} fails on $G_\ell$. Moreover the proportion of the first $n$ ants whose
reinforced path $\gamma_k$ is not a geodesic converges almost surely to $(1-1/\ell)\rho(1-\rho)$, and
\[
\PP(\gamma_n\text{ is not a geodesic})\longrightarrow\bigl(1-\tfrac1\ell\bigr)\,\EE[\rho(1-\rho)]>0 .
\]
\end{corollary}

\begin{corollary}\label{cor:losange}
For the loop-erased ant process on the losange, almost surely $W_e(n)/n\to\chi_e$ for every edge $e$, with
$\chi_{Nu}=\chi_{uF}=1-\chi_{Nv}=1-\chi_{vF}\in(0,1)$ and $\chi_{uv}=0$.
\end{corollary}

Corollary~\ref{cor:losange} is the case $\ell=1$ of Theorem~\ref{thm:main}; it is the statement
of~\cite[Theorem~1.4]{KMS} with the loop-erased rule in place of the uniform-geodesic rule, and it agrees with
Conjecture~\ref{conj}. The smallest graph of the family on which the conjecture fails is $G_2$, which has six vertices
and seven edges. The half of the conjecture which says that the edges of the geodesics keep a positive share of the
weight holds on every $G_\ell$.

The proof rests on two facts. The first is classical: the reinforced path of an ant has the law of the path from $N$
to $F$ in a random spanning tree of the weighted graph (Lemma~\ref{lem:ust}). On $G_\ell$ the spanning trees are easy
to list, and this gives the law of $\gamma_{n+1}$ given $W(n)$ in closed form (Lemma~\ref{lem:law}). The legs enter
through the factor $\ell$, the number of places where a leg can be cut: the total weight of the trees whose path from
$N$ to $F$ uses the rung carries this factor, while the weight of the trees that contain both legs does not, and for
$\ell\ge2$ this is enough to keep the rung alive. The second fact is an exact martingale (Section~\ref{sec:mart}), which gives the
convergence of one coordinate of the process and from which the rest follows by elementary arguments.

\section{The law of the reinforced path}

\begin{lemma}\label{lem:ust}
Let $G$ be a finite connected graph with positive edge weights $(w_e)$ and let $N\ne F$ be two vertices. Let $X$ be
the weighted random walk started at $N$ and stopped when it reaches $F$, and let $\gamma$ be the loop-erasure of the
reversed trajectory. Then $\gamma$ has the law of the path between $F$ and $N$ in a random spanning tree $T$ of $G$
with $\PP(T=t)$ proportional to $\prod_{e\in t}w_e$.
\end{lemma}
\begin{proof}
For a vertex $x\ne N$ visited by the walk let $\sigma_x$ be the first time at which the walk is at $x$, and let
$p(x)=X_{\sigma_x-1}$. In the reversed trajectory the last visit to a vertex $x$ takes place at time $K-\sigma_x$, and
the next vertex is $X_{\sigma_x-1}=p(x)$. By the definition of the loop-erasure, $\gamma$ is therefore the path
$F,p(F),p(p(F)),\dots$, followed until it reaches $N$. Now let the walk continue after time $K$ until it has visited
every vertex. By the theorem of Aldous and Broder~\cite{Aldous,Broder} (see also~\cite[Chapter~4]{LP}), the edges
$\{x,p(x)\}$, $x\ne N$, form a spanning tree $T$ with the law in the statement. The path from $F$ to $N$ in $T$ is
$F,p(F),p(p(F)),\dots$, because the first-visit times decrease along this sequence, and these vertices and their
first-entrance edges are determined by the walk up to time $K$. Hence this path is $\gamma$.
\end{proof}

On $G_\ell$ there are four simple paths from $N$ to $F$,
\begin{align*}
\pi_1&=(N,u,x_1,\dots,F), & \pi_3&=(N,u,v,y_1,\dots,F),\\
\pi_2&=(N,v,y_1,\dots,F), & \pi_4&=(N,v,u,x_1,\dots,F).
\end{align*}
Each of them contains all the edges of one leg and no edge of the other. Hence in the ant process all the edges of the
$u$-leg have the same weight at all times, and likewise for the $v$-leg. We write
\[
A_n=W_{Nu}(n),\quad B_n=W_{Nv}(n),\quad R_n=W_{uv}(n),
\]
and $C_n$, $D_n$ for the common weight of the edges of the $u$-leg and of the $v$-leg. Since every $\pi_i$ contains
exactly one of $Nu$, $Nv$ and exactly one of the two legs,
\begin{equation}\label{eq:sum}
A_n+B_n=C_n+D_n=n+2=:S_n .
\end{equation}

\begin{lemma}\label{lem:law}
Give weight $A$ to $Nu$, $B$ to $Nv$, $R$ to $uv$, $C$ to each edge of the $u$-leg and $D$ to each edge of the $v$-leg
of $G_\ell$. The path from $N$ to $F$ in the random spanning tree of Lemma~\ref{lem:ust} equals $\pi_i$ with
probability $q_i=m_i/(m_1+m_2+m_3+m_4)$, where
\[
m_1=AC(D+\ell B+\ell R),\qquad m_2=BD(C+\ell A+\ell R),\qquad m_3=\ell ARD,\qquad m_4=\ell BRC,
\]
and $m_1+m_2+m_3+m_4=CD(A+B)+\ell AB(C+D)+\ell R(A+B)(C+D)$.
\end{lemma}
\begin{proof}
The graph $G_\ell$ consists of three internally disjoint paths between $u$ and $v$: the rung, the path $u,N,v$, and
the path $u,x_1,\dots,F,\dots,y_1,v$ made of the two legs. It has $2\ell+2$ vertices and $2\ell+3$ edges, and its
spanning trees are obtained by deleting one edge from each of two of the three paths; there are $2+6\ell$ of them. We
list them with their weights divided by $C^{\ell-1}D^{\ell-1}$.

Delete the rung and $Nu$: the path is $\pi_2$ and the weight is $BCD$. Delete the rung and $Nv$: $\pi_1$, weight $ACD$.
Delete the rung and one of the $\ell$ edges of the $u$-leg: $\pi_2$, weight $ABD$ each. Delete the rung and an edge of
the $v$-leg: $\pi_1$, weight $ABC$ each. Delete $Nu$ and an edge of the $u$-leg: $\pi_2$, weight $BRD$ each. Delete
$Nu$ and an edge of the $v$-leg: $\pi_4$, weight $BRC$ each. Delete $Nv$ and an edge of the $u$-leg: $\pi_3$, weight
$ARD$ each. Delete $Nv$ and an edge of the $v$-leg: $\pi_1$, weight $ARC$ each. Summing over the trees with a given
path gives $m_1,\dots,m_4$.
\end{proof}

\section{The drift}

Let $\cF_n$ be the $\sigma$-field generated by the first $n$ ants, and put
\[
a_n=\frac{A_n}{S_n},\quad b_n=\frac{B_n}{S_n}=1-a_n,\quad c_n=\frac{C_n}{S_n},\quad d_n=\frac{D_n}{S_n}=1-c_n,
\quad h_n=\frac{R_n}{S_n} .
\]
By Lemmas~\ref{lem:ust} and~\ref{lem:law} and by~\eqref{eq:sum}, $\PP(\gamma_{n+1}=\pi_i\mid\cF_n)=q_i(n)$ with
\begin{equation}\label{eq:q}
q_1=\frac{ac(\ell b+d+\ell h)}{Q},\quad q_2=\frac{bd(\ell a+c+\ell h)}{Q},\quad q_3=\frac{\ell ahd}{Q},\quad
q_4=\frac{\ell bhc}{Q},\qquad Q=\ell ab+cd+\ell h,
\end{equation}
where $a,b,c,d,h,Q$ are evaluated at time $n$. Note that $0<h_n<1$ and $ab,cd\le\frac14$, so
\begin{equation}\label{eq:Qbound}
\ell\,a_nb_n+c_nd_n\le Q_n\le\tfrac{5\ell+1}{4} .
\end{equation}

\begin{lemma}\label{lem:drift}
With $g=ad+bc-ab-cd/\ell$,
\begin{align*}
\PP(Nu\in\gamma_{n+1}\mid\cF_n)&=q_1+q_3=a+\frac{\ell ab(c-a)}{Q}, &
\PP(Nv\in\gamma_{n+1}\mid\cF_n)&=b-\frac{\ell ab(c-a)}{Q},\\
\PP(ux_1\in\gamma_{n+1}\mid\cF_n)&=q_1+q_4=c+\frac{cd(a-c)}{Q}, &
\PP(vy_1\in\gamma_{n+1}\mid\cF_n)&=d-\frac{cd(a-c)}{Q},\\
\PP(uv\in\gamma_{n+1}\mid\cF_n)&=q_3+q_4=h+\frac{\ell h(g-h)}{Q}, &&
\end{align*}
all quantities on the right being evaluated at time $n$.
\end{lemma}
\begin{proof}
Using $c+d=1$,
$q_1+q_3-a=a\bigl(\ell bc+cd+\ell hc+\ell hd-\ell ab-cd-\ell h\bigr)/Q=\ell ab(c-a)/Q$. Using $a+b=1$,
$q_1+q_4-c=c\bigl(\ell ab+ad+\ell ah+\ell bh-\ell ab-cd-\ell h\bigr)/Q=cd(a-c)/Q$. Finally
$q_3+q_4-h=\ell h\bigl(ad+bc-ab-cd/\ell-h\bigr)/Q$. The two remaining identities follow because $\gamma_{n+1}$
contains exactly one of $Nu$, $Nv$ and exactly one of the two legs.
\end{proof}

When $a=c=p$ one has $g=(1-1/\ell)p(1-p)$. The points $a=c=p$, $h=(1-1/\ell)p(1-p)$, $0<p<1$, are thus the zeros of
the drift with $h>0$ when $\ell\ge2$, and Theorem~\ref{thm:main} says that the process converges to a random one of
them. When $\ell=1$ one has $g=(a-c)^2$.

\section{Proof of Theorem~\ref{thm:main}}

For an integer $k\ge1$ let $\psi(k)=\sum_{j=1}^{k-1}1/j$, so that $\psi(k+1)-\psi(k)=1/k$. Comparing the sum with the
integral of $1/x$,
\begin{equation}\label{eq:psi}
0\le\psi(k)-\log k\le1,\qquad k\ge1,
\end{equation}
and $\psi(k)-\log k$ is non-decreasing in $k$ and converges to Euler's constant. In particular, if $(k_n)$ is a
non-decreasing sequence of positive integers, then $\psi(k_n)-\log k_n$ converges. We also use the following remark
several times: if $(k_n)$ is a non-decreasing sequence of positive integers with $k_{n+1}-k_n\in\{0,1\}$, then
\begin{equation}\label{eq:pi}
\sum_{n\ge0}\frac{k_{n+1}-k_n}{k_n^2}\le\sum_{j\ge1}\frac1{j^2}=\frac{\pi^2}{6},
\end{equation}
because each integer is the value of $k_n$ at no more than one time $n$ with $k_{n+1}=k_n+1$.

\subsection{A martingale}\label{sec:mart}
Let
\[
M_n=\psi(A_n)-\psi(B_n)+\ell\bigl(\psi(C_n)-\psi(D_n)\bigr),\qquad
r_n=\log\frac{A_n}{B_n}+\ell\log\frac{C_n}{D_n}=\logit a_n+\ell\,\logit c_n,
\]
where $\logit x=\log(x/(1-x))$.

\begin{lemma}\label{lem:mart}
$(M_n)$ is a martingale bounded in $L^2$. Consequently $M_n$ converges almost surely and in $L^2$, $r_n$ converges
almost surely to a finite limit $r_\infty$, and $r^*:=\sup_n|r_n|<\infty$ almost surely.
\end{lemma}
\begin{proof}
We have $\psi(A_{n+1})-\psi(A_n)=\one\{Nu\in\gamma_{n+1}\}/A_n$, and similarly for $B$, $C$, $D$. By
Lemma~\ref{lem:drift}, with all quantities at time $n$,
\begin{align*}
\EE\bigl[\psi(A_{n+1})-\psi(A_n)\mid\cF_n\bigr]&=\frac{1}{aS}\Bigl(a+\frac{\ell ab(c-a)}{Q}\Bigr)
=\frac1S\Bigl(1+\frac{\ell b(c-a)}{Q}\Bigr),\\
\EE\bigl[\psi(B_{n+1})-\psi(B_n)\mid\cF_n\bigr]&=\frac1S\Bigl(1-\frac{\ell a(c-a)}{Q}\Bigr),\\
\EE\bigl[\psi(C_{n+1})-\psi(C_n)\mid\cF_n\bigr]&=\frac1S\Bigl(1+\frac{d(a-c)}{Q}\Bigr),\\
\EE\bigl[\psi(D_{n+1})-\psi(D_n)\mid\cF_n\bigr]&=\frac1S\Bigl(1-\frac{c(a-c)}{Q}\Bigr).
\end{align*}
Hence
\[
\EE[M_{n+1}-M_n\mid\cF_n]=\frac{1}{SQ}\Bigl(\ell(a+b)(c-a)+\ell(c+d)(a-c)\Bigr)=0 .
\]
Write $M_{n+1}-M_n=U_n+\ell V_n$, where $U_n$ equals $1/A_n$ or $-1/B_n$ according to whether $\gamma_{n+1}$ contains
$Nu$ or $Nv$, and $V_n$ equals $1/C_n$ or $-1/D_n$ according to the leg. By~\eqref{eq:pi} applied to $A$, $B$, $C$ and
$D$, $\sum_nU_n^2\le\pi^2/3$ and $\sum_nV_n^2\le\pi^2/3$, so that $\sum_n(M_{n+1}-M_n)^2\le2(1+\ell^2)\pi^2/3$ surely.
Since $M_0=0$, it follows that $\EE[M_n^2]\le2(1+\ell^2)\pi^2/3$ for all $n$, and $M_n$ converges almost surely and in
$L^2$.

Next, $M_n-r_n$ is the sum of the four terms $\pm(\psi(A_n)-\log A_n)$, $\dots$, $\pm\ell(\psi(D_n)-\log D_n)$. Each
of the sequences $A_n,B_n,C_n,D_n$ is non-decreasing, so each term converges, and by~\eqref{eq:psi}
$|M_n-r_n|\le1+\ell$. Therefore $r_n$ converges almost surely to a finite limit and $\sup_n|r_n|<\infty$.
\end{proof}

\subsection{The process stays away from the boundary}

\begin{lemma}\label{lem:away}
Almost surely there is a (random) $\kappa_0>0$ such that $a_n,b_n,c_n,d_n\ge\kappa_0$ for all $n\ge0$.
\end{lemma}
\begin{proof}
Let $\kappa=(1+e^{r^*/(1+\ell)})^{-1}$, so that $0<\kappa\le\frac12$. We first note that
\begin{equation}\label{eq:impl}
a_n<\kappa\quad\Longrightarrow\quad a_n\le c_n .
\end{equation}
Indeed, if $a_n>c_n$ then $\logit c_n<\logit a_n$, hence $-r^*\le r_n<(1+\ell)\logit a_n$, which gives $a_n>\kappa$.

Let $Y_n=\psi(A_n)-\log S_n$. By~\eqref{eq:psi}, $Y_n-1\le\log a_n\le Y_n$. Let
\[
\Delta_n=\frac{\one\{Nu\in\gamma_{n+1}\}}{A_n}-\frac1{S_n}\Bigl(1+\frac{\ell b_n(c_n-a_n)}{Q_n}\Bigr),\qquad
\mathcal N_n=\sum_{k<n}\Delta_k .
\]
By the first display in the proof of Lemma~\ref{lem:mart}, $(\mathcal N_n)$ is a martingale. Its increments satisfy
$\EE[\Delta_n^2\mid\cF_n]\le\EE[\one\{Nu\in\gamma_{n+1}\}/A_n^2\mid\cF_n]$, and by~\eqref{eq:pi} the sum over $n$ of the
expectations of the right-hand side is at most $\pi^2/6$. So $\mathcal N_n$ is bounded in $L^2$, it converges almost
surely, and $\mathcal N^*:=\sup_n|\mathcal N_n|<\infty$ almost surely. Since $\log(1+1/S_n)\le1/S_n$,
\begin{equation}\label{eq:Yinc}
Y_{n+1}-Y_n=\frac{\one\{Nu\in\gamma_{n+1}\}}{A_n}-\log\Bigl(1+\frac1{S_n}\Bigr)
\ge\Delta_n+\frac{\ell b_n(c_n-a_n)}{Q_nS_n} .
\end{equation}

Fix $n$ with $a_n<\kappa$, and let $m$ be the last time before $n$ with $a_m\ge\kappa$; it exists because
$a_0=\frac12\ge\kappa$. For $m<k<n$ we have $a_k<\kappa$, so $a_k\le c_k$ by~\eqref{eq:impl} and the last term
of~\eqref{eq:Yinc} is non-negative at time $k$. Summing~\eqref{eq:Yinc} over $m<k<n$ gives
$Y_n\ge Y_{m+1}-2\mathcal N^*$. Moreover $Y_{m+1}\ge Y_m-\log(1+1/S_m)\ge Y_m-\log\frac32$ and
$Y_m\ge\log a_m\ge\log\kappa$. Hence
\[
\log a_n\ge Y_n-1\ge\log\kappa-\log\tfrac32-1-2\mathcal N^* .
\]
The same bound holds trivially when $a_n\ge\kappa$. So $\inf_na_n>0$ almost surely.

The map exchanging $u$ with $v$ and $x_i$ with $y_i$ is an automorphism of $G_\ell$ that fixes $N$ and $F$; it
exchanges $a_n$ with $b_n$ and $c_n$ with $d_n$. Hence $\inf_nb_n>0$ almost surely as well. Finally
$\ell\,\logit c_n=r_n-\logit a_n$ is bounded, so $c_n$ and $d_n$ are bounded away from $0$.
\end{proof}

\subsection{The two geodesics}

\begin{lemma}\label{lem:delta}
Almost surely $a_n-c_n\to0$, and $a_n\to\rho$, $c_n\to\rho$, where $\rho=(1+e^{-r_\infty/(1+\ell)})^{-1}\in(0,1)$.
\end{lemma}
\begin{proof}
Let $\delta_n=a_n-c_n$ and $\xi_{n+1}=\one\{Nu\in\gamma_{n+1}\}-\one\{ux_1\in\gamma_{n+1}\}$. Since
$a_{n+1}-a_n=(\one\{Nu\in\gamma_{n+1}\}-a_n)/(S_n+1)$ and similarly for $c_n$,
\[
\delta_{n+1}-\delta_n=\frac{\xi_{n+1}-\delta_n}{S_n+1},\qquad
\EE[\xi_{n+1}-\delta_n\mid\cF_n]=-\lambda_n\delta_n,\qquad \lambda_n=\frac{\ell a_nb_n+c_nd_n}{Q_n},
\]
by Lemma~\ref{lem:drift}. As $|\xi_{n+1}-\delta_n|\le2$,
\[
\EE[\delta_{n+1}^2\mid\cF_n]\le\delta_n^2-\frac{2\lambda_n\delta_n^2}{S_n+1}+\frac{4}{(S_n+1)^2} .
\]
By the theorem of Robbins and Siegmund~\cite{RS}, $\delta_n^2$ converges almost surely and
$\sum_n\lambda_n\delta_n^2/(S_n+1)<\infty$ almost surely. By Lemma~\ref{lem:away} and~\eqref{eq:Qbound},
$\lambda_n\ge4(\ell+1)\kappa_0^2/(5\ell+1)>0$ for all $n$, and $\sum_n1/(S_n+1)=\infty$; so the limit of $\delta_n^2$
is $0$.

The numbers $a_n$ and $c_n$ stay in $[\kappa_0,1-\kappa_0]$, on which $\logit$ is uniformly continuous; hence
$\logit a_n-\logit c_n\to0$, and with Lemma~\ref{lem:mart}, $(1+\ell)\logit a_n\to r_\infty$. This gives the limit
of $a_n$, and $c_n=a_n-\delta_n$ has the same limit.
\end{proof}

\subsection{The rung}

\begin{lemma}\label{lem:det}
Let $(\beta_n)$ be positive numbers with $\beta_n\to0$ and $\sum_n\beta_n=\infty$, let $(g_n)$ be real numbers
converging to $g^*\ge0$, and let $(\eta_n)$ be a convergent real sequence. Let $(u_n)$ satisfy
\[
u_{n+1}-u_n=\beta_n\bigl(g_n-e^{x_n}\bigr),\qquad x_n=u_n+\eta_n,
\]
and assume that $|g_n-e^{x_n}|\le2$ for all $n$. Then $e^{x_n}\to g^*$.
\end{lemma}
\begin{proof}
Fix $\theta>g^*$ and $\varepsilon>0$, and choose $n_1$ such that for all $m,n\ge n_1$ we have $g_n\le(g^*+\theta)/2$,
$2\beta_n\le\varepsilon$ and $|\eta_m-\eta_n|\le\varepsilon$. If $n\ge n_1$ and $x_n>\log\theta$, then
$u_{n+1}-u_n\le-\beta_n(\theta-g^*)/2$. If this held for all large $n$, then $u_n\to-\infty$ and, $\eta_n$ being
bounded, $x_n\to-\infty$, which is absurd. So there is $k\ge n_1$ with $x_k\le\log\theta$. Let $n>k$ with
$x_n>\log\theta$ and let $m$ be the last time before $n$ with $x_m\le\log\theta$, so that $m\ge k$. For $m<j<n$ we
have $u_{j+1}\le u_j$, hence $u_n\le u_{m+1}\le u_m+2\beta_m\le u_m+\varepsilon$ and
$x_n\le x_m+2\varepsilon\le\log\theta+2\varepsilon$. Therefore $\limsup_nx_n\le\log\theta+2\varepsilon$, and since
$\theta>g^*$ and $\varepsilon>0$ are arbitrary, $\limsup_ne^{x_n}\le g^*$. This proves the lemma when $g^*=0$.

If $g^*>0$, fix $0<\theta<g^*$ and $\varepsilon>0$, and choose $n_1$ such that for all $m,n\ge n_1$ we have
$g_n\ge(g^*+\theta)/2$, $2\beta_n\le\varepsilon$ and $|\eta_m-\eta_n|\le\varepsilon$. If $n\ge n_1$ and
$x_n<\log\theta$, then $u_{n+1}-u_n\ge\beta_n(g^*-\theta)/2$; this cannot hold for all large $n$, since then
$x_n\to+\infty$. The same last-exit argument, with the inequalities reversed, gives
$\liminf_nx_n\ge\log\theta-2\varepsilon$, and so $\liminf_ne^{x_n}\ge g^*$.
\end{proof}

\begin{lemma}\label{lem:rung}
Almost surely $h_n\to(1-1/\ell)\rho(1-\rho)$.
\end{lemma}
\begin{proof}
Let $g_n=a_nd_n+b_nc_n-a_nb_n-c_nd_n/\ell$. By Lemma~\ref{lem:delta}, $g_n\to g^*:=(1-1/\ell)\rho(1-\rho)\ge0$.
Let $Y'_n=\psi(R_n)-\log S_n$. By Lemma~\ref{lem:drift},
\[
\EE\bigl[\psi(R_{n+1})-\psi(R_n)\mid\cF_n\bigr]=\frac{1}{h_nS_n}\Bigl(h_n+\frac{\ell h_n(g_n-h_n)}{Q_n}\Bigr)
=\frac1{S_n}+\beta_n(g_n-h_n),\qquad\beta_n=\frac{\ell}{Q_nS_n} .
\]
Let $\Delta'_n=\one\{uv\in\gamma_{n+1}\}/R_n-1/S_n-\beta_n(g_n-h_n)$ and $\mathcal N'_n=\sum_{k<n}\Delta'_k$. As in
the proof of Lemma~\ref{lem:away}, $(\mathcal N'_n)$ is a martingale bounded in $L^2$, so it converges almost surely.
With $e_n=1/S_n-\log(1+1/S_n)\in[0,1/(2S_n^2)]$ we get
\[
Y'_{n+1}-Y'_n=\Delta'_n+\beta_n(g_n-h_n)+e_n .
\]
Put $\varepsilon_n=\psi(R_n)-\log R_n$, which converges because $R_n$ is non-decreasing, and
\[
u_n=Y'_0+\sum_{k<n}\beta_k(g_k-h_k),\qquad\eta_n=\mathcal N'_n+\sum_{k<n}e_k-\varepsilon_n .
\]
Then $\eta_n$ converges almost surely, $\log h_n=Y'_n-\varepsilon_n=u_n+\eta_n$, and
$u_{n+1}-u_n=\beta_n(g_n-e^{u_n+\eta_n})$. We have $-\frac12\le g_n\le1$ and $0<h_n<1$, so $|g_n-h_n|\le2$. By
Lemma~\ref{lem:away} and~\eqref{eq:Qbound}, $(\ell+1)\kappa_0^2\le Q_n\le(5\ell+1)/4$, so $\beta_n\to0$ and
$\sum_n\beta_n=\infty$. Lemma~\ref{lem:det} gives $h_n\to g^*$.
\end{proof}

\subsection{End of the proof}
Since $S_n/n\to1$, Lemmas~\ref{lem:delta} and~\ref{lem:rung} give the limits in Theorem~\ref{thm:main}. The
automorphism used in the proof of Lemma~\ref{lem:away} exchanges $\rho$ and $1-\rho$, so they have the same law. By
Lemma~\ref{lem:away} the weights $A_n,B_n,C_n,D_n$ tend to infinity, so the four terms of $M_n-r_n$ in the proof of
Lemma~\ref{lem:mart} tend to $\pm$ Euler's constant and $\pm\ell$ times Euler's constant and cancel; thus
$r_\infty=\lim_nM_n$. As $M_n$ converges in $L^2$, $\EE[r_\infty]=0$ and
$\EE[r_\infty^2]\ge\EE[(M_1-M_0)^2]>0$, since $M_1-M_0=1+\ell$ when $\gamma_1=\pi_1$. So $r_\infty$, and with it
$\rho$, is not almost surely constant. This proves Theorem~\ref{thm:main}.

\begin{proof}[Proof of Corollary~\ref{cor:fail}]
The paths $\pi_3$ and $\pi_4$ are the simple paths from $N$ to $F$ that are not geodesics, and they are the ones that
contain the rung. So the number of $k\le n$ for which $\gamma_k$ is not a geodesic is $R_n-1$, and
$(R_n-1)/n\to(1-1/\ell)\rho(1-\rho)$ almost surely, which is positive when $\ell\ge2$. By Lemma~\ref{lem:drift},
$\PP(\gamma_{n+1}\in\{\pi_3,\pi_4\}\mid\cF_n)=h_n+\ell h_n(g_n-h_n)/Q_n$, which converges almost surely to the same
limit because $Q_n$ is bounded below; the last statement follows by bounded convergence.
\end{proof}

\section{Simulations}

The program \texttt{sim\_ants.c} runs the process as defined in the introduction: it stores the walk of each ant,
reverses it, erases the loops in chronological order and adds one to the weights of the edges of the resulting path.
It makes no use of Lemma~\ref{lem:ust} or Lemma~\ref{lem:law}.

With the weights held fixed, the frequencies of the four paths agree with Lemma~\ref{lem:law}. For $\ell=2$ and
$(A,B,C,D,R)=(2,3,1,4,2)$, two million ants gave the frequencies $0.15544$, $0.60037$, $0.17762$ and $0.06657$ for
$\pi_1,\dots,\pi_4$, and Lemma~\ref{lem:law} gives $0.15556$, $0.60000$, $0.17778$ and $0.06667$. For $\ell=3$ and
the weights $(5,2,3,4,1)$ the frequencies were $0.44208$, $0.38081$, $0.13629$ and $0.04083$, and the lemma gives
$0.44218$, $0.38095$, $0.13605$ and $0.04082$. The script
\texttt{check\_path\_law.py} verifies Lemma~\ref{lem:law} symbolically for $\ell\le5$ by listing the spanning trees,
and verifies the identities of Lemma~\ref{lem:drift}.

Table~\ref{tab:runs} shows the state of the process after $10^8$ ants in independent runs. In every run $a_n$ and
$c_n$ differ by less than $0.0003$, and for $\ell\ge2$ the value of $h_n$ is close to $(1-1/\ell)a_n(1-a_n)$.

\begin{table}[ht]
\centering
\begin{tabular}{rrcccc}
\toprule
$\ell$ & seed & $a_n$ & $c_n$ & $h_n$ & $(1-1/\ell)a_n(1-a_n)$\\
\midrule
2 & 1 & 0.38569 & 0.38574 & 0.12126 & 0.11847\\
2 & 2 & 0.94939 & 0.94938 & 0.02419 & 0.02402\\
2 & 3 & 0.14083 & 0.14079 & 0.06048 & 0.06050\\
2 & 4 & 0.31259 & 0.31262 & 0.11194 & 0.10744\\
2 & 5 & 0.53892 & 0.53897 & 0.12692 & 0.12424\\
2 & 6 & 0.12526 & 0.12523 & 0.05685 & 0.05478\\
3 & 1 & 0.64536 & 0.64537 & 0.15316 & 0.15258\\
3 & 2 & 0.93299 & 0.93305 & 0.04243 & 0.04168\\
5 & 1 & 0.53274 & 0.53295 & 0.19958 & 0.19914\\
5 & 2 & 0.41410 & 0.41408 & 0.19433 & 0.19410\\
10 & 1 & 0.35721 & 0.35711 & 0.20658 & 0.20665\\
1 & 1 & 0.46134 & 0.46132 & 0.02736 & 0\\
1 & 2 & 0.91163 & 0.91162 & 0.00766 & 0\\
\bottomrule
\end{tabular}
\caption{The loop-erased ant process on $G_\ell$ after $n=10^8$ ants.}\label{tab:runs}
\end{table}

The convergence of $h_n$ is slow when $\ell$ is small. Linearising the drift of $h$ at the limit point gives the rate
$-(1-1/\ell)/2$, which suggests that $h_n-(1-1/\ell)a_n(1-a_n)$ is of order $n^{-(1-1/\ell)/2}$, that is $n^{-1/4}$
for $\ell=2$. In the run with $\ell=2$ and seed $1$ this difference was $0.0401$, $0.0187$, $0.0090$, $0.0050$,
$0.0028$ at $n=10^4,\dots,10^8$; the last ratios are $0.56$, and $10^{-1/4}=0.562$. For $\ell=1$ the drift of $h$ is
$h(\delta^2-h)/Q$ with $\delta=a-c$, so that $1/h_n$ should grow by about $\log10/Q_n$ per decade of $n$. In the run with $\ell=1$ and seed $1$, $h_n$ was
$0.0519$, $0.0423$, $0.0357$, $0.0311$, $0.0274$ at $n=10^4,\dots,10^8$, so $1/h_n$ grew by between $4.1$ and $4.4$
per decade, and $\log10/Q_n$ is close to $4.4$ there.

\section{Remarks}

\begin{remark}
On $G_\ell$ the limit of the rung is a function of the limit $\rho$ of the two geodesics, and all edges of the
geodesics survive. We do not know which edges survive on a general graph, nor whether the edges of the geodesics
always do. By Lemma~\ref{lem:ust} the mean drift of the process on any graph is given by the laws of the paths
between $N$ and $F$ in weighted random spanning trees, and on $G_\ell$ the rung survives because of the weight
$CD(A+B)$ of the two trees that contain both legs, which is not multiplied by $\ell$ in Lemma~\ref{lem:law}: at
$a=c=p$ and small $h$ the rung is reinforced at $2\ell/(\ell+1)$ times the rate $h$ at which it is diluted, and this
ratio exceeds $1$ exactly when $\ell\ge2$.
\end{remark}

\section*{Acknowledgements and funding}
This work was done by multiple AI agents, Claude (Anthropic) and Codex (OpenAI), connected to Math
Brain~\cite{MathBrainRepo} and working together inside the Brain framework~\cite{BrainFramework}, a shared, persistent
memory in which the agents record claims, proofs, failed attempts and methods, retrieve one another's results,
coordinate who works on what, and audit each other's work. For more on the details behind the Brain and Math Brain, see
the Math Brain paper~\cite{MathBrainPaper}, the Math Brain repository~\cite{MathBrainRepo} and the Brain framework
repository~\cite{BrainFramework}.

The research presented in this paper was supported by the European Research Council (ERC) under the European Union's
Horizon Europe research and innovation programme (grant agreement No.~101041711), by the Simons Foundation as part of
the Collaboration on the Mathematical and Scientific Foundations of Deep Learning, by Heights Labs, by Convex Nexus
Capital, by the Israel Science Foundation (grant number 2258/19), and by the Israel Science Foundation (ISF Grant
4101/25). Computational resources were provided in part by Google Cloud through its education credits program (credit
530550103).

\section*{Declaration of competing interest}
The author holds equity in Heights Labs and in Convex Nexus Capital and has a management role at Convex Nexus Capital.
These entities are included in the funding declaration above.

\section*{Declaration of generative AI and AI-assisted technologies}
The paper comes out of Math Brain, an ongoing project on the use of AI systems for research mathematics, and its
results were developed with the Brain~\cite{BrainFramework} and Math Brain~\cite{MathBrainRepo}. In October 2026 an AI
agent connected to the Brain, working from the project's list of open problems in probability, found the graph $G_2$
and showed that on it the rung keeps a positive share of the weight with positive probability; a Codex agent (OpenAI)
and a Claude agent (Anthropic) each checked that argument. In November 2026 a Claude Fable agent (Anthropic) found
Theorem~\ref{thm:main} and its proof, wrote the programs, ran the simulations and prepared the manuscript. These
checks are not peer review. The author directed the work and is responsible for the content of the paper and for
decisions concerning its dissemination.

\section*{Code and data availability}
Programs and outputs:\\
\url{https://github.com/yspennstate/loop-erased-ants-non-geodesic-edge}.

The paper is part of the Math Brain project, whose public repository is
\url{https://github.com/yspennstate/math-brain-public}~\cite{MathBrainRepo}.

\begin{thebibliography}{99}
\bibitem{Aldous} D.~J.~Aldous, The random walk construction of uniform spanning trees and uniform labelled trees,
\emph{SIAM J. Discrete Math.} \textbf{3} (1990), 450--465.
\bibitem{Aletti} G.~Aletti, A.~Collevecchio, I.~Crimaldi, M.~Deijfen, S.~Dereich, A.~Ghiglietti, R.~van der Hofstad,
M.~Holmes, D.~Kious, V.~Kleptsyn, B.~Lodewijks, P.~M\"orters and B.~Schapira, Open problems in stochastic
reinforcement models, report of the MATRIX--MFO Tandem Workshop \emph{Stochastic Reinforcement Processes and Graphs},
March 2023.
\bibitem{Broder} A.~Broder, Generating random spanning trees, in \emph{30th Annual Symposium on Foundations of Computer
Science}, IEEE, 1989, 442--447.
\bibitem{KMS} D.~Kious, C.~Mailler and B.~Schapira, Finding geodesics on graphs using reinforcement learning,
\emph{Ann. Appl. Probab.} \textbf{32} (2022), no.~5; arXiv:2010.04820.
\bibitem{KMS2} D.~Kious, C.~Mailler and B.~Schapira, The trace-reinforced ants process does not find shortest paths,
\emph{J. \'Ec. polytech. Math.} \textbf{9} (2022), 505--536.
\bibitem{LP} R.~Lyons and Y.~Peres, \emph{Probability on Trees and Networks}, Cambridge University Press, 2016.
\bibitem{MV} C.~Mailler and Z.~Varin, The two-nest ants process on triangle-series-parallel graphs,
arXiv:2601.22855.
\bibitem{RS} H.~Robbins and D.~Siegmund, A convergence theorem for non negative almost supermartingales and some
applications, in \emph{Optimizing Methods in Statistics}, Academic Press, 1971, 233--257.
\bibitem{MathBrainRepo} Y.~Shmalo, \emph{Math Brain: public snapshot, code and tests}, 2026,
\url{https://github.com/yspennstate/math-brain-public}.
\bibitem{MathBrainPaper} Y.~Shmalo, \emph{Math Brain: shared, auditable memory for multi-agent mathematical research},
October 2026, \url{https://github.com/yspennstate/math-brain-public/tree/main/papers/math-brain-2026-10-24}.
\bibitem{BrainFramework} Y.~Shmalo, \emph{Brain Framework}, public software repository, 2026,
\url{https://github.com/yspennstate/brain-framework-public}.
\end{thebibliography}
\end{document}
